% ============================================================
% SECTION 2 — Imam al-Hassan's Pacifism (Chapter 2, Sessions 1-2)
% ============================================================
\section{al-Hassan's Pacifism}

\begin{frame}{Not a Question of Temperament}
\bigquote{Different temperaments cannot explain why two brothers, raised the same way, chose opposite methods.}{Mutahhari's opening move, paraphrased}
\begin{itemize}[itemsep=5pt]
  \item The lazy answer: al-Hassan was peaceful, al-Husayn was brave
  \item Mutahhari's answer: both were reasoned applications of the \emph{same} jurisprudence to \emph{different} circumstances
  \item Proof by precedent: even the Prophet alternated war and peace (Hudaybiyyah vs. later battles) depending on circumstance, not temperament
\end{itemize}
\end{frame}

\begin{frame}{Categories of Jihad}
\centering\scriptsize
\renewcommand{\arraystretch}{1.4}
\begin{tabular}{|p{3.1cm}|p{5.4cm}|p{2.6cm}|}
\hline
\textbf{Category} & \textbf{Condition} & \textbf{Applies to} \\\hline
Offensive jihad & Requires an imam's presence; not permitted to a lone authority acting alone & Neither brother's case \\\hline
Defensive jihad & Obligatory on all when Islamic land/people/property attacked & Closer to al-Husayn \\\hline
Jihad-equivalent (self-defense) & Fighting to save one's own life; scales with what's at stake & al-Husayn at Karbala \\\hline
Fighting Muslim rebels (Qur'an 49:9) & Reconcile first; fight only the side that refuses & Available, not obligatory, to al-Hassan \\\hline
Legitimate truce & Only if too weak to prevail, need time, or truce may draw the other side toward Islam; fixed term only & al-Hassan's treaty \\\hline
\end{tabular}
\vspace{0.3cm}
\delivernote{Walk the room through this table before saying a word about the actual treaty -- that ordering is deliberate, so nobody judges the treaty by gut feeling before knowing the legal categories.}
\end{frame}

\begin{frame}{Hudaybiyyah: The Direct Precedent}
\begin{itemize}[itemsep=5pt]
  \item Apparently one-sided 10-year truce: Muslim defectors returned, no reciprocal clause
  \item 'Umar objected; the Prophet accepted anyway
  \item Within under two years: more conversions than the previous two decades of open conflict
\end{itemize}
\bigquote{Leave me and seek someone else... I would rather stay a minister than to become a chief.}{'Ali -- a parallel restraint, decades earlier}
\end{frame}

\begin{frame}{Five Contrasts: al-Hassan vs. al-Husayn}
\centering\scriptsize
\renewcommand{\arraystretch}{1.4}
\begin{tabular}{|l|p{4.7cm}|p{4.7cm}|}\hline
&\textbf{al-Hassan}&\textbf{al-Husayn}\\\hline
Position & Sitting caliph -- death parallels 'Uthman's assassination & No office -- death is honor, not institutional collapse\\\hline
Military balance & Iraq (30k--100k max) vs. Mu'awiyah's 150k; Kufah exhausted, bribed & Different balance a generation later\\\hline
Allegiance demand & None -- treaty explicitly excludes any oath & Yazid's explicit personal demand -- the trigger\\\hline
Kufah's mood & Exhausted, hostile -- wore armor even to lead prayer & 18{,}000 letters -- a mass mandate history would fault him for ignoring\\\hline
Ruler's record & Mu'awiyah's tyranny not yet proven; offered signed blank paper & Two decades of documented injustice by Yazid's time\\\hline
\end{tabular}
\delivernote{No single row is decisive alone -- the strength of the argument is that on every axis Mutahhari can find, the two situations genuinely diverge.}
\end{frame}

\begin{frame}{What the Treaty Actually Secured}
\begin{itemize}[itemsep=4pt]
  \item Governance to Mu'awiyah \emph{conditioned} on ruling by Qur'an and Sunnah
  \item Succession reverts to al-Hassan, then al-Husayn -- explicitly temporary
  \item End to public cursing of 'Ali in Syria
  \item Treasury funds for al-Hassan's household and Siffin/Jamal war-widows
  \item General amnesty for 'Ali's former supporters
\end{itemize}
\bigquote{}{Mu'awiyah later broke essentially every one of these terms -- which is what exposed his true character, not the peace itself.}
\end{frame}

\qaframe{al-Hassan's Pacifism}{%
  \qa{Lay out all five contrasts. Does any single one, alone, justify al-Hassan's peace -- or does the argument only work cumulatively?}
     {No single factor is decisive alone. Removing any one wouldn't flip the verdict, since the others still stand independently. The strength is that on every axis Mutahhari can find, the two brothers' situations genuinely diverge.}
  \qa{If Hudaybiyyah is the model, what is the risk of treating every costly-looking peace as automatically wise in hindsight?}
     {Survivorship bias -- Hudaybiyyah is remembered as vindicated because it worked. The specific jurisprudential conditions (weakness, need for time, likelihood of drawing the other side toward Islam) have to actually hold in a given case, not just be assumed by loose analogy to a famous precedent.}
  \qa{Play devil's advocate: what is the strongest counter-argument a critic of al-Hassan could still make?}
     {The five contrasts explain why war would have been strategically unwise, but don't fully settle whether al-Hassan had a moral duty to resist regardless of the odds. Mutahhari's implicit answer: obligation scales with genuine possibility and manifest cause, not a blanket principle that surrender is always shameful.}
}
